\chapter{Witt Vector Algebra}\label{appendix:witt_vector_algebra}

This appendix collects the general algebra of finite-length \(p\)-typical Witt
vectors used in Chapter~2. The chapter itself retains the
Artin--Schreier--Witt classification and ramification theory; here we record the
Witt-polynomial construction, the ghost-map arguments, and the structural
identities needed there.

\providecommand{\bW}{\mathbb{W}}

\section{The ring of Witt vectors}

For $n \geq 0$, the $n$-th \textbf{Witt polynomial} in the variables $X_0, \dots, X_n$
is
\[
w_n(X_0, \dots, X_n) \;=\; \sum_{j=0}^{n} p^j X_j^{p^{n-j}}
\;=\; X_0^{p^n} + p X_1^{p^{n-1}} + \dots + p^n X_n .
\]

\begin{theorem}[Witt; {\cite[II, \S6, Theorem~6]{serreLF}}, {\cite[\S1]{Rabinoff2014Witt}}]
\label{theorem:witt_ring_structure}
There exists a unique family of polynomials
$S_n, P_n \in \Z[X_0, \dots, X_n; Y_0, \dots, Y_n]$, $n \geq 0$, such that
\[
w_n(S_0, \dots, S_n) = w_n(X_0,\dots,X_n) + w_n(Y_0,\dots,Y_n),
\qquad
w_n(P_0, \dots, P_n) = w_n(X_0,\dots,X_n) \cdot w_n(Y_0,\dots,Y_n)
\]
for all $n$. For any commutative ring $A$, the operations
\[
\vec a \oplus \vec b := \bigl(S_n(\vec a; \vec b)\bigr)_{0 \le n \le r-1},
\qquad
\vec a \odot \vec b := \bigl(P_n(\vec a; \vec b)\bigr)_{0 \le n \le r-1}
\]
make the set $A^r$ into a commutative ring $W_r(A)$, the ring of \textbf{$p$-typical
Witt vectors of length $r$ over $A$}, with zero $(0, \dots, 0)$ and unit
$(1, 0, \dots, 0)$. Additive inverses are likewise given by universal integral
polynomials; we write $\ominus$ for Witt subtraction. The \textbf{ghost map}
\[
w \colon W_r(A) \longrightarrow A^r,
\qquad
\vec a \longmapsto \bigl(w_0(\vec a), w_1(\vec a), \dots, w_{r-1}(\vec a)\bigr)
\]
is a ring homomorphism to the product ring $A^r$; it is injective when $p$ is a
non-zero-divisor in $A$, and bijective when $p \in A^\times$.
\end{theorem}

We index the components of a Witt vector $\vec a = (a_0, \dots, a_{r-1})$ by
$0, \dots, r-1$. For $r = 1$ the polynomials reduce to $S_0 = X_0 + Y_0$ and
$P_0 = X_0 Y_0$, so $W_1(A) = A$ as a ring.

Since the ring operations are given by universal polynomials with integer
coefficients, $W_r$ is a functor: a ring homomorphism $\sigma \colon A \to B$ induces
the ring homomorphism
\[
W_r(\sigma) \colon W_r(A) \to W_r(B),
\qquad
W_r(\sigma)(a_0, \dots, a_{r-1}) = (\sigma a_0, \dots, \sigma a_{r-1}),
\]
acting \emph{componentwise}. Two consequences of this trivial-looking observation are
used repeatedly and deserve emphasis: a ring automorphism $\sigma$ of $A$ commutes
with $\oplus$ and $\ominus$ on $W_r(A)$, and a Witt vector $\vec a \in W_r(A)$
satisfies $W_r(\sigma)(\vec a) = \vec a$ if and only if \emph{each component} $a_n$ is
$\sigma$-invariant. In particular, for a group $\Gamma$ acting on $A$ by ring
automorphisms,
\begin{equation}\label{eq:witt_invariants_componentwise}
W_r(A)^{\Gamma} \;=\; W_r\bigl(A^{\Gamma}\bigr)
\qquad \text{inside } W_r(A).
\end{equation}

The uniqueness statement in \Cref{theorem:witt_ring_structure} is an instance of a
general principle that we invoke several times, so we record it explicitly.

\begin{lemma}[ghost principle]\label{lemma:ghost_principle}
Let $\Phi$ and $\Psi$ be two operations on Witt vectors given by universal polynomials
with integer coefficients in the components of their arguments. If
$w \circ \Phi = w \circ \Psi$ holds over the polynomial ring
$\Z[X_0, X_1, \dots]$ on the components of the arguments, then $\Phi = \Psi$ over
every commutative ring.
\end{lemma}

\begin{proof}
Over $\Z[X_0, X_1, \dots]$ the prime $p$ is a non-zero-divisor, so the ghost map is
injective by \Cref{theorem:witt_ring_structure} and $\Phi = \Psi$ there; this is an
identity of integral polynomials, which persists under every base change
$\Z[X_0, X_1, \dots] \to A$.
\end{proof}

The only structural feature of the addition polynomials $S_n$ that we shall need is
their ``triangular'' shape:

\begin{prop}[shape of Witt addition]\label{prop:witt_addition_shape}
For every $n \geq 0$ there is a polynomial
$C_n \in \Z[X_0, \dots, X_{n-1}; Y_0, \dots, Y_{n-1}]$, involving only the variables of
index \emph{strictly less} than $n$, such that
\begin{equation}\label{eq:witt_addition_shape}
(\vec a \oplus \vec b)_n \;=\; a_n + b_n + C_n(a_0, \dots, a_{n-1}; b_0, \dots, b_{n-1}).
\end{equation}
We refer to $C_n$ as the \textbf{carry polynomial}. The analogous statement holds for
$\ominus$.
\end{prop}

\begin{proof}
By induction on $n$, using the recursion defining $S_n$:
\[
p^n S_n \;=\; \sum_{j \leq n} p^j \bigl(X_j^{p^{n-j}} + Y_j^{p^{n-j}}\bigr)
\;-\; \sum_{j < n} p^j S_j^{p^{n-j}} .
\]
The terms with $j = n$ contribute $p^n (X_n + Y_n)$, and all remaining terms involve
only variables and polynomials $S_j$ of index $j < n$, which by induction lie in
$\Z[X_{<n}; Y_{<n}]$. Since $S_n$ has integer coefficients
(\Cref{theorem:witt_ring_structure}), $C_n := S_n - X_n - Y_n$ is an integral
polynomial in the variables of index $< n$. The statement for $\ominus$ follows by the
same argument applied to the inversion polynomials.
\end{proof}

\section{Frobenius, Verschiebung and truncation}

From now on, $A$ denotes an $\F_p$-algebra; this is the only case we use. Three
operators act on the rings $W_r(A)$:

\begin{definition}\label{definition:witt_operators}
Let $A$ be an $\F_p$-algebra.
\begin{enumerate}
    \item The \textbf{Frobenius} $F \colon W_r(A) \to W_r(A)$ is
    $F := W_r(\varphi_A)$, where $\varphi_A \colon a \mapsto a^p$ is the absolute
    Frobenius of $A$. Explicitly, $F$ acts componentwise:
    $F(a_0, \dots, a_{r-1}) = (a_0^p, \dots, a_{r-1}^p)$. By functoriality, $F$ is a
    ring endomorphism of $W_r(A)$.
    \item The \textbf{Verschiebung} is the shift
    \[
    V \colon W_{r-1}(A) \to W_r(A),
    \qquad
    V(a_0, \dots, a_{r-2}) = (0, a_0, \dots, a_{r-2}).
    \]
    Iterating, $V^j \colon W_{r-j}(A) \to W_r(A)$; in particular
    $V^{r-1} \colon A = W_1(A) \to W_r(A)$ sends $c \mapsto (0, \dots, 0, c)$.
    \item The \textbf{truncation} is
    \[
    t \colon W_r(A) \to W_{r-1}(A),
    \qquad
    t(a_0, \dots, a_{r-1}) = (a_0, \dots, a_{r-2}).
    \]
\end{enumerate}
\end{definition}

\begin{rem*}
For a general (not necessarily $\F_p$-) algebra $A$, the Frobenius is defined by
universal polynomials characterized by $w_n \circ F = w_{n+1}$, and the componentwise
formula above holds precisely in characteristic $p$; see
\cite[II, \S6]{serreLF} or \cite[\S2]{Rabinoff2014Witt}. We will not need the general
case.
\end{rem*}

We record from the general theory (\cite[II, \S6]{serreLF},
\cite[\S2]{Rabinoff2014Witt}) how multiplication by $p$ interacts with these
operators.

\begin{prop}\label{prop:witt_multiplication_by_p}
Let $A$ be an $\F_p$-algebra. Then:
\begin{enumerate}
    \item $t$ is a natural ring homomorphism commuting with $F$, and its kernel is
    \[
    \ker\bigl(t \colon W_r(A) \to W_{r-1}(A)\bigr)
    \;=\; V^{r-1} A \;=\; \{(0, \dots, 0, c) : c \in A\}.
    \]
    \item $V$ is additive, and $F V = V F$.
    \item Multiplication by $p$ on $W_r(A)$ equals $V \circ F = F \circ V$
    (composed with the appropriate truncation); explicitly,
    \[
    p \cdot (a_0, \dots, a_{r-1}) \;=\; (0,\, a_0^p,\, \dots,\, a_{r-2}^p).
    \]
    Iterating, $p^j \cdot \vec a = (0, \dots, 0, a_0^{p^j}, \dots, a_{r-1-j}^{p^j})$;
    in particular $p^r = 0$ on $W_r(A)$, and
    \begin{equation}\label{eq:witt_p_to_r_minus_1}
    p^{r-1} \cdot (a_0, \dots, a_{r-1}) \;=\; \bigl(0, \dots, 0,\, a_0^{p^{r-1}}\bigr).
    \end{equation}
\end{enumerate}
\end{prop}

\begin{proof}
(1) That $t$ is a ring homomorphism follows from the recursive definition of the
$S_n, P_n$: the first $r-1$ components of a sum or product depend only on the first
$r-1$ components of the arguments (\Cref{prop:witt_addition_shape}). Naturality and
the commutation with the (componentwise) $F$ are clear, and the kernel is computed
componentwise. (2) Additivity of $V$ follows from \Cref{lemma:ghost_principle} and
the ghost identity $w_n(V\vec a) = p\, w_{n-1}(\vec a)$ (with $w_{-1} := 0$); the
commutation $FV = VF$ is immediate since $F$ is componentwise. (3) The identity
$F V = p$ holds over arbitrary rings for the general Frobenius, by
\Cref{lemma:ghost_principle} applied to
$w_n(FV\vec a) = w_{n+1}(V \vec a) = p\, w_n(\vec a) = w_n(p \vec a)$; in
characteristic $p$ the general Frobenius is componentwise, giving the displayed
formula. See \cite[II, \S6, Proposition~8 ff.]{serreLF} for details.
\end{proof}

\begin{example}\label{example:witt_of_Fp}
$W_r(\F_p) \cong \Z/p^r\Z$ canonically. Indeed, by
\Cref{eq:witt_p_to_r_minus_1} the unit $\underline 1 = (1, 0, \dots, 0)$ satisfies
$p^j \cdot \underline 1 = (0, \dots, 0, 1, 0, \dots, 0)$ ($1$ in position $j$), so
$\underline 1$ has additive order $p^r$; since $|W_r(\F_p)| = p^r$, the additive group
is cyclic of order $p^r$ generated by the unit, and $m \mapsto m \cdot \underline 1$
is a ring isomorphism $\Z/p^r\Z \xrightarrow{\;\cong\;} W_r(\F_p)$. We write
$\underline m \in W_r(\F_p)$ for the image of $m \in \Z/p^r\Z$. (Similarly, the full
ring of Witt vectors satisfies $W(\F_p) \cong \Z_p$; we shall not need this.)
\end{example}

\section{Isobaricity of the Witt addition polynomials}

The following homogeneity property of the universal addition polynomials is what
makes degree estimates on Witt sums possible: even though the carry polynomials $C_n$
of \Cref{eq:witt_addition_shape} are complicated, they are \emph{isobaric} for the
weighting in which the $j$-th component has weight $p^j$.

\begin{lemma}[isobaricity of Witt addition]\label{lemma:witt_addition_isobaric}
Assign to the variables $X_j$ and $Y_j$ the weight $p^j$. Then the $n$-th addition
polynomial $S_n$ is isobaric of weight $p^n$: every monomial of $S_n$ has total weight
exactly $p^n$. The same holds for each component of a $d$-fold iterated Witt sum
$\bigoplus_{i=1}^{d} \vec a_i$, viewed as a universal polynomial in the components
$a_{i,j}$ of weight $p^j$.
\end{lemma}

\begin{proof}
The $n$-th ghost polynomial $w_n = \sum_{j \le n} p^j X_j^{p^{n-j}}$ is isobaric of
weight $p^n$ for this weighting. The recursion
\[
p^n S_n = \sum_{j \le n} p^j \bigl(X_j^{p^{n-j}} + Y_j^{p^{n-j}}\bigr)
- \sum_{j < n} p^j S_j^{p^{n-j}}
\]
preserves isobaricity of weight $p^n$ by induction on $n$: each $S_j$ is isobaric of
weight $p^j$, so $S_j^{p^{n-j}}$ is isobaric of weight $p^n$. Isobaricity persists
under iterated substitution of Witt sums into Witt sums (substituting an isobaric
polynomial of weight $p^j$ for a variable of weight $p^j$ preserves weights), giving
the statement for $d$-fold sums.
\end{proof}
